\section{Dataset}
\label{sec:data}

\textbf{Radio.} We use the F1 Team Radio dataset~\cite{ciletti2025f1radio}: 14{,}681 broadcast clips (44.8\,h of audio; 43 driver channels; 149 Grands Prix; 2018--2025) with ASR transcripts. The 2022 season is absent and 2025 is thin (669 clips over 23 races, versus $\sim$130 per race in earlier seasons). We removed 204 byte-identical clips that recur across races (all transcribed ``Thank you.''; most contain no speech) and 107 clips whose transcript is contained in a longer clip from the same driver and race, leaving 14{,}370 clips.

\textbf{Telemetry.} OpenF1~\cite{openf1} has no data before 2023, so we use FastF1~\cite{fastf1}, which reads the same official timing feed for every season. Laps, car data (speed, throttle, binary brake, RPM, gear, DRS at $\sim$3.7\,Hz), race-control messages and weather were obtained for 148 of 149 races (the 2018 Italian GP is missing from the source; 2018 Bahrain has laps but no car data).

\textbf{Join and timestamp delay.} A clip's \texttt{message\_timestamp} is in UTC (the time included in clip IDs is local track time, which accounts for hour-long differences). Associating each clip with its car's lap positions 72\% of clips during the race, 15\% post-chequered flag, and 13\% pre-start. We assessed the delay between the broadcast time and the speech using two anchors. (i) "Box" calls are vocalized before entering the pit, but 81.6\% of 163 calls are timestamped \emph{after} the pit-in of the car (median $+31.5$\,s). (ii) The first clip after the race for each car appears at least $\sim$64\,s following the finish line (5th percentile). Consequently, timestamps trail spoken content by approximately 30–60 seconds, roughly one lap, leading to the use of lap-level windows throughout.


\section{Methodology}
\label{sec:method}

The pipeline consists of eight predetermined stages (A--H), each concluded by a checkpoint where the outcome is evaluated against a preset threshold~\cite{nosek2018prereg}. Decisions regarding analysis were documented prior to the creation of the relevant labels (\texttt{PREREGISTRATION*.md} in the repository); any deviations are recorded with timestamps and presented here.


\subsection{Evaluation Splits}
All divisions are organized by race: each segment of a Grand Prix is placed in one group, as segments from a single race have common track, weather, and events. We employ (i) 5-fold race-grouped cross-validation (CV) as the main method due to the limited chronological test season; (ii) a chronological split (train 2018--2023, validate 2024, test 2025); and (iii) leave-one-driver-out (LODO) for the eight most common drivers, \emph{crossed} with the race folds ensuring that each held-out fold remains separate from training by both driver and race. A unit test fails if any race occurs in two instances.


\subsection{Stage A: Speaker Attribution}
A channel carries the driver's \emph{and} the engineer's voice, usually in one recording without turn markers. We diarize every clip with pyannote's \texttt{speaker-diarization-community-1} pipeline~\cite{bredin2023pyannote,plaquet2023powerset} (batch size 4, which keeps peak GPU memory at 0.25\,GB rather than 10.9\,GB on our 8\,GB card and runs at 49$\times$ real time; output identical to the default, DER $=0$ on 40 clips), giving 29{,}134 turns of at least 0.5\,s. Within each channel with $\geq$25 clips, the pipeline's WeSpeaker~\cite{wang2023wespeaker} speaker embeddings are clustered (cosine, average linkage, $k\in\{2..6\}$ by silhouette); each cluster is named \textsc{driver} or \textsc{engineer} by the duration-weighted mean of rule-based linguistic cues (imperatives, ``box'', lap targets, the driver's first name as a vocative $\rightarrow$ engineer; first-person experience of the car $\rightarrow$ driver), with a constraint that each channel has at least one cluster of each role. Small channels use a logistic regression on embeddings trained on the large channels' cluster names; turns whose cluster centroid similarity is in the channel's bottom decile fall back to the text cues.

\textbf{Evaluation.} One annotator labelled 360 randomly drawn turns (stratified by season, blind to predictions and transcripts) plus 20 repeats for test--retest agreement. Accuracy is reported with a race-bootstrap 95\% CI~\cite{efron1979bootstrap}; the pre-registered gate was 0.85.

\subsection{Stage B: Re-Transcription}
Whisper large-v3~\cite{radford2023whisper} (beam 5, word timestamps, no conditioning on previous text) was run on all clips; it is \emph{not} fine-tuned. We adapt it through prompting only: a season-specific \texttt{initial\_prompt} lists that season's driver names (from official results) and core F1 jargon. Word error rate (WER) against 41 hand-transcribed clips (859 words; drawn before any ASR output existed, stratified by season) compares the dataset's transcripts, plain Whisper and prompted Whisper, with a paired bootstrap CI on relative improvement. The pre-registered bar for adoption was a 15\% relative WER reduction.

\subsection{Stage C: Annotation}
A claim taxonomy was defined backwards from the available sensors: \textsc{grip\_deg}, \textsc{traffic}, \textsc{rel\_pace}, \textsc{incident}, \textsc{power}, \textsc{no\_sensor} (claims public telemetry cannot check, e.g.\ brake temperature or fuel), plus \textsc{no\_claim} for acknowledgements and questions, which the original taxonomy lacked. Arousal was labelled on a 3-point scale from the voice. We use \emph{arousal} rather than ``stress'': vocal arousal here is confounded with g-load, breathing and broadcast compression~\cite{russell1980circumplex}. One annotator labelled 450 utterances drawn from turns attributed to the driver (80\% in-race; equal per season), marking those that were not the driver.

\subsection{Stage D: Telemetry and the Lap-Time Model}
An expected-lap-time model is fit per race by least squares on clean laps (excluding lap 1, pit laps, safety-car/VSC/red-flag laps and laps $>$107\% of the driver's median): a driver intercept, lap number (fuel-burn proxy), tyre age per compound, and a traffic indicator (car ahead within 1.5\,s). The \emph{residual} (actual minus expected) is the outcome for both T2 and T4, so its own error is reported with 5-fold within-race CV.

\subsection{Stage E: Models}
\textbf{T1/T3 text: full fine-tuning with a shared trunk.} DeBERTa-v3-base~\cite{he2023debertav3} is fully fine-tuned with a shared encoder and two heads (claim type and arousal), trained with the sum of two class-weighted cross-entropies (lr $2\times10^{-5}$, 8 epochs, batch 16, max length 64; all fixed in advance). Baselines: majority class; a \emph{context-only} logistic regression on lap, position, tyre age and gap, with no text; TF-IDF logistic regression; and an \emph{LLM zero-shot} baseline. For the latter, a small open instruction-tuned model (Qwen2.5-1.5B-Instruct~\cite{qwen2025qwen25}, run locally) receives the category definitions from the annotation guide and the utterance; each candidate label is scored by the model's log-likelihood of that label as the answer, softmaxed over candidates so that AUC-PR is computed exactly as for trained models. One prompt was written and run once, with no prompt tuning.

\textbf{T3 audio: partial fine-tuning with a learned layer weighting.} WavLM-base-plus~\cite{chen2022wavlm} has its convolutional encoder and lower six transformer layers frozen; the upper six layers are fine-tuned (lr $2\times10^{-5}$) together with a softmax-normalised learned weight per upper layer~\cite{yang2021superb} and a linear arousal head that also receives season indicators (the broadcast codec changed over the years). As an ablation, a frozen model uses a fixed mean of layers 4--9 plus prosodic features (energy, spectral centroid, zero-crossing rate, F0), z-scored within driver-race, with logistic regression; a fused model stacks audio features with out-of-fold text probabilities.

\textbf{T2 telemetry consistency.} We call the target \emph{telemetry-consistency}, never truth: apart from FIA decisions on incidents, no veracity ground truth exists. (a) \emph{Rules.} Before any claim label existed we pre-registered, per claim type, a lap-level rule with a primary threshold and a sensitivity sweep (e.g.\ \textsc{grip\_deg}: mean clean-lap residual over laps $L{-}2..L \geq 0.3$\,s; \textsc{traffic}: gap ahead $<1.5$\,s or a blue flag; \textsc{incident}: a race-control message naming the car in $L{-}1..L{+}2$). Each rule is evaluated for every lap of every car (161{,}914 car-laps), so that a claim's consistency can be compared with the car's base rate, with other claims, and with a placebo. (b) \emph{Neural.} Because a model trained on rule outputs could only reproduce the rules, the neural T2 instead predicts the \emph{human} claim label from telemetry alone: a 1\,Hz car-data window (120\,s before to 60\,s after the timestamp, skewed earlier because of the measured lag; z-scored within driver-session) passes through a 1-D CNN, and seven lap tokens ($L{-}4..L{+}2$: residual, gaps, position, tyre age, maximum speed and RPM, pit/neutralised/race-control flags) through a linear layer. Each candidate claim type is an embedding that queries the telemetry tokens by multi-head cross-attention~\cite{vaswani2017attention}; a sigmoid gate mixes the attended evidence with the query (constrained fusion). Following evidence that naive concatenation of text into time-series models underperforms constrained fusion~\cite{lee2026cfa}, we ablate this against concatenating mean-pooled telemetry with the claim embedding.

\subsection{Stage F: Controls and Metrics}
(i) \emph{Reverse-time placebo}: every T2 analysis is repeated on the window three laps \emph{before} the message. Equal performance backwards means the signal predates the claim, which indicates the broadcaster's selection of eventful moments rather than the driver's report. (ii) The context-only baseline. (iii) A leakage audit: the race-grouping test, and a null in which 30 sets of random labels give mean macro ROC-AUC $0.499$ (SD 0.035).

\textbf{Metrics.} Per-class AUC-PR (average precision~\cite{davis2006prroc,saito2015pr}) with 95\% CIs from resampling races. Because AUC-PR of a random ranking exceeds prevalence when positives are few, every table reports a simulated chance level (mean and 95th percentile over 200 random rankings of the same labels).

\subsection{Stage G: Arousal and Subsequent Performance (T4)}
For each in-race lap $L$ with a driver utterance, $x$ is the mean predicted probability of raised or high arousal: soft probabilities rather than hard labels, since measurement error attenuates coefficients~\cite{fuller1987measurement}. Outcomes are the mean clean-lap residual over $L{+}1..L{+}3$ and the residual SD over $L{+}1..L{+}5$, controlling for the residual at lap $L$ (drivers may be aroused \emph{because} the lap was slow), tyre age, race progress and position. A linear mixed model~\cite{seabold2010statsmodels} with driver random intercepts and random slopes on $x$ is compared with an intercept-only model by a likelihood-ratio test using the boundary mixture $\tfrac12\chi^2_1+\tfrac12\chi^2_2$~\cite{self1987lrt}; we report the slope variance component, not per-driver rankings. The minimum detectable effect (MDE; 80\% power, $\alpha=0.05$, with a Kish design effect~\cite{kish1965survey} for clustering by driver) is computed from the controls-only model and written to disk \emph{before} the arousal model is fit~\cite{simmons2011false}. An attenuation check re-estimates the effect on the hand-labelled subset using labels instead of probabilities.

